
\begin{personalcv}{攻读学位期间公开发表论文}[addtotoc]

首先，列出在攻读硕士学位期间发表与学位论文有关的学术论文（含已录用），并注明属于学位论文内容的部分（章节），作者（最多三个）、论文题目、刊物名称、时间、卷期号、页码以及检索信息、与学位论文相关章节。{\bf\itshape 在攻读硕士学位期间以外的时间或与学位论文内容(章节)无关的论文不得列出。}

其次，列出在攻读硕士学位期间参加学术会议发表的会议论文，参与的科研项目（如国家自然科学基金或国家“863”计划等），发明专利、科研奖励等。

书写格式说明：

本人的姓名应加粗和标记下划线。每段落首行缩进2字（即：将每个成果作为一个独立的段落）。

例：

1. 发表学术论文
\begin{enumerate}[nosep, label=$\lbrack$\arabic*$\rbrack$]
  \item{\underline{\bf Xiaoming Zhang}, Yi Li, John R. E., et al. Carbon isotope evidence for the stepwise oxidation of the Proterozoic environment [J]. Nature, 1992,359(1):605-609. (SCI检索号：123DX) （本学位论文第一章）.}
\end{enumerate}{}\vspace{1em}

2. 会议论文
\begin{enumerate}[nosep, label=$\lbrack$\arabic*$\rbrack$]
  \item{\underline{\bf 张晓明}，会议论文题目，会议名称，口头报告或墙报，会议地点，时间。}
  \item{\underline{\bf Xiaoming Zhang}, Development of cellular biology. Proceedings of the Fifth Canadian Mathematical Congress (oral presentation), Tokyo, 2018.}
\end{enumerate}{}\vspace{1em}

3. 参与科研项目
\begin{enumerate}[nosep, label=$\lbrack$\arabic*$\rbrack$]
  \item{国家自然科学基金项目(51276055)：西南喀斯特山区土地利用和土地覆被变化及其对土地资源可持续性影响研究，2013.1 – 2016.12，负责人：李文。} %负责人一般应为指导教师
\end{enumerate}{}\vspace{1em}

4. 发明专利
\begin{enumerate}[nosep, label=$\lbrack$\arabic*$\rbrack$]
  \item{\underline{\bf 张晓明}，发明人2，发明人3. 多功能一次性压舌板:中国,92214985.2[P]. 发明类别：发明专利，公开（或授权）日期：1993,04,14。}
\end{enumerate}{}\vspace{1em}

5. 获得奖励情况
\begin{enumerate}[nosep, label=$\lbrack$\arabic*$\rbrack$]
  \item{“大型C/E复合材料构件高质高效加工关键技术及其工艺装备”，机械工业科学技术奖-科技进步一等奖，2013.10，本人排序第1。}
\end{enumerate}

\end{personalcv}


